\section{Critical estimates and their geometric consequences}\label{sec:inputs}

The proof uses critical estimates at two different levels.  Uniform arm
estimates, valid down to the mesh, control the expansion in sites.  At fixed
positive scales, the scalar Cardy hypothesis also gives quenched crossing
limits, and these permit more precise planar constructions.  We establish
both levels here.  In particular, the estimates for exploration near a
boundary will be proved in conformal coordinates; a Euclidean half-plane
estimate in the physical plane would not suffice for the domains produced
by a stopped exploration.

\subsection{Arm estimates, planar conventions, and localization}

We suppress the model superscript when a statement applies to both models.
The parameter $e$ is the mesh, and the number of sites per unit area is
\[
 m_e=e^{-2}\quad\text{in the Poisson model},\qquad
 m_e=\frac{2}{\sqrt3e^2}\quad\text{on the triangular lattice}.
\]
Fix a sufficiently large absolute constant $B_0$ and put $b_e=B_0e$.
Write $A_e(r,R)$ for the annealed probability of four alternating arms
across a square annulus with inner and outer radii $r,R$.  The center is a
lattice site in the triangular model.  Bounded changes of radii, and the
replacement of square by circular annuli with margins, will be absorbed
by quasi-multiplicativity.  Arms of the same color are required to use
disjoint sites, including the equal-colored neighboring arms in an odd
arm sequence.

\begin{lemma}[Uniform critical arm inputs]\label{lem:arm-inputs}
For each model there are constants $C,c,\alpha,\beta>0$ and
$g\in(1,2)$ such that, for $b_e\le r\le s\le R$,
\begin{align}
 c(r/R)^g&\le A_e(r,R)\le C(r/R)^{1+\alpha},\label{eq:four-powers}\\
 C^{-1}A_e(r,s)A_e(s,R)&\le A_e(r,R)
       \le C A_e(r,s)A_e(s,R).\label{eq:qm}
\end{align}
The one-arm probability is at most $C(r/R)^\beta$, and the
polychromatic five-arm probability is at most $C(r/R)^2$.
The alternating three-arm probability in a half-plane in a standard
orientation is at most $C(r/R)^2$.

In the Poisson model, if $X_e(\eta;r,R)$ is the conditional probability
of any of these specified arm events, then
\begin{equation}\label{eq:arm-second-moment}
 \bigl(\EE X_e(\eta;r,R)^2\bigr)^{1/2}
       \le C\,\PP(\text{the corresponding arm event}).
\end{equation}
At fixed positive radii with $r\le R/2$, the conditional arm probabilities
minus their annealed probabilities converge to zero in probability over
the geometry.  Both assertions apply to the half-plane events just stated.
For triangular four arms one may choose $g\in(5/4,4/3)$.
\end{lemma}

These are critical inputs from the literature, with their precise scope
as follows.  For Poisson--Voronoi percolation, rescaling by $e^{-1}$ gives
unit intensity, independent fair colors, and inner radii at least one.
The full-plane second-moment comparison and concentration are
\cite[Theorem~1.4]{vanneuville-quantitative}; the power estimates and
quasi-multiplicativity are supplied by
\cite[Propositions~3.1--3.4 and~6.2]{vanneuville-quantitative}.
The half-plane second-moment statement uses Remark~4.1 and equation~(3.2)
there.  For its concentration assertion it is enough to use the weaker
variance estimate proved in Section~4, equation~(4.1).  At unit intensity,
this bounds the variance divided by the square of the annealed arm
probability by $Cr^{-h}$, for some $h>0$, $r\ge r_0$, and $r\le R/2$.
After rescaling, this relative variance is at most $C(e/r)^h$, with the
same half-plane modification.  No half-plane version of the sharper
variance formula in Theorem~1.4 is needed.

The crossing and gluing method originates in the work of Russo and
Seymour--Welsh~\cite{russo78,seymourwelsh78}.
For the triangular lattice we use critical RSW, arm separation and
quasi-multiplicativity as in \cite[Sections~2.3 and~4]{nolin08}, together with
the arm exponents of \cite{smirnovwerner01}; see also
\cite[Theorems~21--22]{nolin08}.  The two-scale inequalities with a
strict exponent margin follow by iterating a sufficiently large fixed
scale ratio and using quasi-multiplicativity.  The finitely many remaining
lattice scales change only the constant.  The estimate for oblique
straight boundaries needed below will be derived after
Lemma~\ref{lem:boundary-arms}; it is not being included implicitly in
the standard-orientation input.

We use an embedded plane triangulation whose vertices are the sites.
For a Voronoi tessellation, draw an edge between neighboring nuclei
through the relative interior of their common cell side.  Its faces may
be curvilinear triangles.  Almost surely the Poisson diagram is in
general position.  The same is true after finitely many fixed nuclei
are inserted, outside a Lebesgue null set of insertion locations.
All arguments take place in bounded regions and therefore involve
locally finite diagrams.  An interface, or \emph{wall}, crosses each
triangle through the edges whose endpoints have different colors.
It can equally be drawn on cell boundaries without changing its order.

Site paths and paths in unions of cells can be interchanged with an
error bounded by the sizes of the cells they use.  In the interior of
the plane, a same-colored turn at a triple vertex can be made through
relative interiors of common sides.  Thus a path with fixed positive
clearance from the sides of a test can be drawn with clearance inside
its color before being trimmed to that test.  For alternating arms,
the intervening opposite-colored paths keep the two paths of either
color separate when this conversion is made in a buffered annulus.

We shall also approximate simple curves by graph paths.  Here is the
elementary approximation fact we use.  In a thin open tube about a
simple arc, follow the successive triangles met by the arc, take paths
on their edges, and erase loops.  Work first in a rectangular coordinate
strip about the arc and let its width decrease more slowly than the
mesh.  The resulting paths converge in curve order: longitudinal
backtracks tend to zero.  For a closed curve, use an annular strip and
retain the simple boundary facing the compact set to be enclosed.
For crosscuts ending on a graph boundary, the same construction works
up to their endpoints.  A boundary cycle can first be made free of
interior chords by retaining the part surrounding a fixed inner disk.
Its inner row is connected along successive vertex stars, so crosscuts
can reach the boundary without a bounded-degree assumption.  Hex
duality in these triangulated disks uses terminal sides, or uncolored
fan terminals attached outside them.

\begin{lemma}[Necessary arms for a local change]\label{lem:pivotal-arms}
Compare two colored tilings which agree outside a ball and throughout
an intervening annulus.  If their crossing indicators for a polygon
quad differ, an annulus contained in the quad and surrounding the
change has four alternating arms.  Near one straight side, at scales
larger than the distance to that side and smaller than the distance
to all vertices and other sides, the corresponding half-annulus has
three alternating arms.  At a vertex or a marked point, a one-arm
test from a neighborhood of that point to a fixed larger distance is
always necessary.  Each assertion allows fixed-factor margins in its
inner and outer radii and uses only colors in the agreement region.
Other prescribed local changes may be present, provided that they
are held fixed in the comparison.
\end{lemma}

\begin{proof}
Use an actual black closed-cell crossing in the configuration with a
crossing.  In the other configuration, closed-set planar duality gives
an open white crossing in the complementary direction, relative to
the quad.  To see the applicable convention, map the quad to a
rectangle and apply the separation theorem for a closed subset and
its open complement.  The white path can be chosen polygonal and
with clearance from the black closed set.  Both paths must visit the
region where the configurations differ.

Trim each path at surrounding curves, using its first hit of the
outer curve and its last preceding hit of the inner one.  Its two
parts lead to its two target sides.  Outside the change, the two
black parts cannot connect to each other: that would give the black
crossing without the change.  The same argument applies to the white
parts, using their open clearance.  Planar separation puts the four
parts in alternating order.  This proves the bulk implication.

Near one straight side, the parts leading to the other three sides
give the three half-plane arms, with the middle one leading to the
opposite side.  The passage to disjoint site paths also respects the
literal closed-cell convention.  A convex cell intersects a straight
side in a connected set; a same-colored shared-cell shortcut either
removes a repeated portion or would already connect the target sides
without the change.  For white arms use relative open paths and pass
through cell interiors and common sides.  Black paths of zero width
on the straight boundary are represented by the cells containing
them, with the allowed margins.  At a corner, at least one target
arc of the crossing is at a fixed positive distance; following the
corresponding part of a crossing supplies the stated one-arm test.
Nothing in this argument changes if other prescribed regions are
held fixed: only the paths in the agreement annulus are used.
\end{proof}

For Poisson geometry, agreement of the nuclei in a neighborhood alone
does not immediately imply agreement of the cells there.  The next
lemma supplies the required screens and also explains how a failed
screen is paid for in an arm estimate.

\begin{lemma}[Buffered localization]\label{lem:localization}
In a shell of radius $t\ge b_e$, one can choose a geometry event
$\mathcal L_e(t)$, determined by the nuclei in a fixed enlargement of
that shell, with
\begin{equation}\label{eq:screen-failure}
 \PP(\mathcal L_e(t)^c)\le C\exp[-c(t/e)^2].
\end{equation}
On this event the tiling and the vertex stars in a smaller shell are
determined there, and every cell meeting that smaller shell has
diameter and distance to its nucleus smaller than any prescribed
fixed fraction of $t$.  Fixed nuclei outside the enlarged shell do
not affect this assertion.  Screens with disjoint enlargements are
independent; arm tests restricted by these screens and their color
bits can likewise be used independently.

Suppose a comparison requires four arms wherever the screens pass
between radii $r,R$.  Its necessary probe, allowing screen failures,
has probability at most $C A_e(r,R)$.  The part in which a screen
fails at a scale $t\ge Mb_e$ is at most
\begin{equation}\label{eq:screen-relative-tail}
 \delta(M)A_e(r,R),\qquad \delta(M)\longrightarrow0,
\end{equation}
when the smaller starting scale is $r\asymp b_e$; the constants may
depend on the bounded factors implicit here.  For a fixed positive
starting radius $r$, the relative failure contribution tends to zero
as $e\to0$.  These conclusions also hold for probes consisting of
four arms up to the distance from a straight side and three
half-plane arms thereafter, once the half-plane estimate is bounded
by $C A_e$.
\end{lemma}

\begin{proof}
Cover the enlarged shell by finitely many squares of side $c_1t$,
where $c_1$ is smaller than the required margins, and require a
nucleus in each square.  The Poisson empty-square formula and a
union bound give \eqref{eq:screen-failure}.  A point in the smaller
shell then has a nucleus within $C c_1t$.  A nucleus beyond the
enlargement cannot be its closest nucleus.  The same comparison
along a segment in a convex cell bounds the entire part of any
cell meeting the smaller shell and hence bounds its star.  This is
also why even a very long cell belonging to a remote prescribed
insertion cannot cross a screened shell.

Use dyadic shells, with spare shells for the buffers.  On complete
success an ordinary arm event with slightly changed radii bounds
the necessary probe.  Otherwise group outcomes by the outermost
failed screen, at scale $t$.  Discard all demands up to scale $Ct$
and keep the arm test beyond that scale, independently of the
failed occupancy test.  Quasi-multiplicativity and
\eqref{eq:four-powers} give, with the convention that an inverted
annulus costs one,
\[
 \exp[-c(t/e)^2]A_e(Ct,R)
 \le C A_e(r,R)\,(t/r)^g\exp[-c(t/e)^2].
\]
The dyadic sum is bounded uniformly.  Its tail with $t\ge Mb_e$
goes to zero after division by $A_e(r,R)$ when $r\asymp b_e$.
When $r$ is fixed, the whole sum goes to zero.  A failure at the
outermost scale is covered by the same inequality using the power
lower bound for $A_e(r,R)$.  This argument does not assert
independence for events involving unrestricted Voronoi cells.

For a side probe, use bulk annuli up to the distance from the side
and then half-annuli centered at the projection onto the side.
Leave disjoint buffers where the two types meet.  The assumed
comparison of the half-plane cost with $A_e$, followed by
quasi-multiplicativity, gives the same calculation.  We shall verify
the comparison in every triangular orientation below.
\end{proof}

\subsection{Quenched crossing limits and deterministic sequences}

Write $\operatorname{Cardy}(Q)=F(x_Q)$ for the scalar value in
Hypothesis~\ref{hyp:cardy}.  The scalar Cardy hypothesis concerns
an average over the geometry.
We next remove that averaging for each fixed polygon test.  The
argument uses the squared arm estimates; ordinary annealed arm
upper bounds alone would not yield the required variance decay.

\begin{proposition}[Quenched scalar Cardy convergence]\label{prop:quenched-cardy}
For every fixed polygon quad $Q$,
\[
 \PP^{\eta_e}(C_Q)\longrightarrow\operatorname{Cardy}(Q)
 \quad\text{in probability over the Poisson geometry}.
\]
The assertion holds simultaneously for every member of any fixed
finite collection of polygon quads.
\end{proposition}

\begin{proof}
Let $Z_e=\PP^{\eta_e}(C_Q)$.  Restrict the relevant geometry to a
large fixed box surrounding $Q$.  The resulting error tends to
zero exponentially as $e\to0$ by localization.  Partition this
box into squares of side $s$, and resample the Poisson process in
one square at a time.  Denote the resulting conditional crossing
probability by $Z_e^{(i)}$.  Efron--Stein gives
\begin{equation}\label{eq:cardy-efron-stein}
 \Var Z_e\le \frac12\sum_i\EE (Z_e-Z_e^{(i)})^2+o_e(1).
\end{equation}
Here and throughout this proof $e\to0$ first, with $s>0$ fixed.
Couple the colors outside the resampled square.  Screens make the
two tilings agree outside a constant enlargement, except on an
event tending to zero.  A difference then requires the arm tests
of Lemma~\ref{lem:pivotal-arms} starting at radius $Cs$.

For squares a fixed distance from all vertices and marks, bulk and
side tests, together with \eqref{eq:arm-second-moment}, contribute
at most $C s^{2+2\alpha}$ per square.  If necessary decrease
$\alpha>0$ so that the side test has at least this exponent.
There are $O(s^{-2})$ such squares.

We spell out the corner summation.  Suppose the distance to the
nearest vertex or mark, capped at a small fixed constant, is
comparable to $u\gg s$.  Probe from $Cs$ to $c_Q u$, using four
arms followed, when necessary, by half-plane arms.  Choose $c_Q$
small enough that the probe meets at most one side, also at narrow
polygon angles.  Independently, with a buffer, require one arm
from $C_Q u$ to a fixed radius whenever the square is near a
corner.  The root-second-moment bound for the resulting necessary
event is
\[
 C(s/u)^{1+\alpha}u^{\beta_0}
\]
for some $\beta_0>0$.  There are $O((u/s)^2)$ squares in this
dyadic distance class.  Their total contribution to
\eqref{eq:cardy-efron-stein} is at most
$C(s/u)^{2\alpha}u^{2\beta_0}$.  The bounded number of squares
within distance $Cs$ of a corner contribute $C s^{2\beta_0}$
by the one-arm test alone.  It follows that
\[
 \limsup_{e\to0}\Var Z_e
 \le C s^{2\alpha}+C s^{2\beta_0}
   +C\sum_{\substack{u\text{ dyadic}\\s\le u\le u_0}}
            (s/u)^{2\alpha}u^{2\beta_0}
 \longrightarrow0
\]
as $s\downarrow0$.  The dyadic sum tends to zero whether
$\alpha$ is smaller than, equal to, or larger than $\beta_0$;
in the equal case the extra factor is only $\log(1/s)$.
The supplied scalar Cardy theorem identifies $\EE Z_e$ and
therefore proves the proposition.  A finite union gives the last
assertion.
\end{proof}

\begin{lemma}[Deterministic sequences and Jordan tests]\label{lem:good-sequences}
From every sequence $e\downarrow0$ in the Poisson model one can
extract a subsequence and realize its geometries on a common
probability space such that, almost surely, the following hold.
The mesh tends to zero on every compact set; crossing probabilities
of all rational polygon quads tend to their Cardy values; and the
fixed-positive-scale arm estimates in Lemma~\ref{lem:arm-inputs}
hold for the conditional color laws, with a fixed constant slack,
on any specified countable family of centers and radii.  Conditional
arm probabilities can also be compared to their mesh-dependent
annealed values at these fixed scales.

Any deterministic sequence satisfying these properties has the
Cardy limit for Jordan quad approximations whose boundaries and
marks converge in curve order to a Jordan quad.  This includes
graph-cycle approximations with fair interior vertices and terminal
sides, and the literal closed-cell crossing convention.  It also
has uniformly positive crossing probabilities along any fixed
finite collection of open tracks with compatible colors.  The
triangular sequence has all these properties deterministically.
\end{lemma}

\begin{proof}
Include the buffered enlargements used below in the countable annulus
family before extracting.
The first assertion is the subsequence principle for convergence
in probability, followed by a diagonal extraction over compact
sets, rational tests, and the specified annuli.  Occupancy gives
the compact mesh assertion; Proposition~\ref{prop:quenched-cardy}
and arm concentration give the others.  At fixed radii the
annealed probabilities are bounded away from zero by RSW and
quasi-multiplicativity.  For $r\le R/2$, absolute concentration therefore
gives the relative comparisons.  For $r>R/2$, the arm event from $r$ to
$2R$ is contained in the event from $r$ to $R$.  Use this inclusion for
the lower comparison and the bound $1$ for the upper comparison;
fixed-ratio RSW bounds the annealed probabilities above and below by
positive constants.  Nearby centers or radii are handled by inclusions
with fixed-factor margins.

Here is the crossing-convention argument.  A Schoenflies strip
about each side of the limiting Jordan quad supplies inner and
outer comparison quads.  For a lower bound, trim in from the two
non-target sides, also near their ends, and extend the target caps
beyond the actual boundary.  Keep positive clearance from the
non-target sides.  Approximate the comparison quad by a rational
polygon.  Any crossing of this polygon gives a crossing of the
original test after trimming at its target sides.  First convert
the crossing to a path through cell interiors and common sides,
using the clearance, and only then trim it.  This order handles
even a crossing initially supplied as a closed connected set.

The complementary color gives the opposite inequality by the
same construction.  As the strip widths decrease, both Cardy
values tend to the desired value, by continuity of conformal maps
under Jordan boundary convergence.  For an approximating graph
disk the comparison paths are trimmed on first meeting its
boundary.  Thus they use only fair interior values and reach
neighbors of the terminal sides; prescribed exterior boundary
colors are not used as parts of the crossing paths.  Outer
comparison tests similarly bound crossings originally expressed
using closed cells.  This proves the claimed sandwich for all
the stated conventions.

For a fixed open track, use a finite chain of overlapping thin
rectangles, together with transverse crossings in their overlaps.
Their crossings concatenate, and such chains can also make
circuits.  Cardy or RSW gives a positive lower bound for each
rectangle.  Harris' inequality~\cite{harris60} gives the product lower bound for
requirements of one color; opposite-colored requirements are
placed in disjoint tracks and use disjoint sites once the mesh is
small.  For the triangular model start with Smirnov's theorem
\cite{smirnov01} for the usual lattice approximations and apply
exactly the same sandwich.  Virtual terminal edges incur only a
vanishing mesh error before the positive-clearance trimming.
\end{proof}

In the rest of this section, a \emph{deterministic sequence} means
one satisfying Lemma~\ref{lem:good-sequences}.  Probabilities then
refer only to its independent fair colors.  An assertion at fixed
scales always means a limit superior along this sequence with
those scales fixed first.  This convention permits us to keep the
Poisson geometry fixed throughout an exploration.

\subsection{Color order and bulk estimates at fixed scales}

The cited alternating arm estimates have to be strengthened before
they can control arbitrary near contacts of paths.  The required
color-order argument is a finite-graph statement and therefore
applies on every one of our deterministic tilings.

\begin{lemma}[Color order in a rectangle]\label{lem:rectangle-colors}
Let a finite plane graph be drawn in a topological rectangle, with
independent fair colors on its vertices.  The probability of $k$
vertex-disjoint crossings between two opposite sides, with a
specified ordered list of colors, does not depend on that list.
The assertion remains valid after arbitrary vertex deletions.
Uncolored fan terminals on the two target sides may be shared;
disjointness and order are then imposed on the paths and their
incident edges off the fan terminals.
\end{lemma}

\begin{proof}
Among the crossings of a prescribed first color, take the leftmost
simple path.  One way to make this choice precise is to add an
external closing edge on the right.  The union of all available
crossing paths with that edge is a union of cycles sharing the
edge and is two-connected.  Its exterior envelope supplies the
leftmost crossing.  Equivalently, it has no bypass to its left
between distinct path vertices, including the terminal incidences.
Whether a proposed path is this leftmost path is determined by
the colors on it and to its left; vertices strictly to its right
remain independent fair bits.

If an ordered family exists, its first path may be replaced by
this leftmost path while all remaining paths stay strictly to
the right.  Conditional on the leftmost path and the revealed
left part, apply induction to the remaining graph.  The
conditional probability of its remaining $k-1$ ordered crossings
does not depend on their colors.  Integrating proves independence
of the last $k-1$ colors.  Global color reversal then removes
dependence on the first one as well.  Vertex deletions change
neither the envelope construction nor the product law on the
unrevealed vertices.
\end{proof}

\begin{lemma}[Bulk arms of arbitrary color order]\label{lem:color-arms}
There are $c_0,C>0$ such that, on a deterministic sequence and at
fixed positive radii $r<R$, the probability of four disjoint bulk
arms of any specified color order has limit superior at most
$C(r/R)^{1+c_0}$.  For six disjoint arms the corresponding bound
is $C(r/R)^{2+c_0}$.  The events may be taken between graph cycles
in buffered circular or square annuli.
\end{lemma}

\begin{proof}
We first compare a color sequence containing both colors with an
alternating sequence.  Choose two adjacent arms of opposite
colors, trim them to proper crossings, and consider the gap
between them.  The disagreement interfaces in its triangulation
contain a strand joining the two annular ends: the switch parity
on the two ends, followed by pairing the interface ends, gives
such a strand.  Its visited triangles have no vertices on the
other arms.  Slitting along it and deleting the vertices of those
triangles leaves a graph in the slit rectangle, possibly with
further vertex deletions, containing the other arms.  Each bank
of the strand already supplies a monochromatic
arm, of opposite colors on its two sides.

For precision, cut the actual disagreement curves at \emph{every}
contact with either annular boundary and sum over their maximal
through components and their two color directions.  Sum also
over the at most $k$ cyclic positions of the opposite-color gap;
this fixes the residual ordered color list and costs only a
factor depending on $k$.  A candidate strand's occurrence
is determined by the bits of its incident triangles.  Conditional
on those bits, the vertices left in the slit rectangle are still
independent and fair.  Lemma~\ref{lem:rectangle-colors} replaces
the ordered colors of the remaining arms by those completing an
alternating family, without changing the conditional probability.
The strand banks lie at the extremes relative to the remaining
paths, so the replaced family and the banks really do give the
alternating arm event.

The overcount in this sum has uniformly bounded moments.  To see
this, let $N$ be the number of through strands after cutting at
boundary contacts.  In a smaller annulus, a bank of a wall gives
a connected site walk: successive triangles connect their
same-side vertices, with short spokes to the wall.  The walk
cannot cross another wall.  After proper trimming it gives an
arm in that gap.  A gap is used by at most its two bordering
strands, so $N$ through strands imply at least a constant times
$N$ disjoint black arms in a buffered subannulus.  If $R/r$ is
larger than a fixed constant, keep a fixed-ratio middle annulus.
A white circuit there has probability at least $c>0$ by
Lemma~\ref{lem:good-sequences}; hence a black crossing has
probability at most $1-c$.  BK on the fixed graph gives an
exponential tail for the number of disjoint black crossings.
Thus $\limsup\EE N^p\le C_p$ for every finite $p$,
uniformly in $R/r$.  Bounded ratios require only an enlarged
constant.

Writing $A_k^{\rm alt}$ for the alternating event in the
same annulus, with the fixed end margins just used absorbed in
its radii, put $P_k^{\rm alt}=\PP(A_k^{\rm alt})$.
For each candidate strand, conditional color-order invariance
replaces its residual-arm event by the completing alternating
event with equal probability.  Summing these replaced
probabilities counts at most $C_kN$ strands on configurations
with an alternating family.  Therefore
\[
 \PP(\text{specified mixed $k$ arms})
 \le C_k\EE[N\mathbf1_{A_k^{\rm alt}}].
\]
The moment bound and H\"older's inequality give
\[
 \limsup\PP(\text{specified mixed $k$ arms})
 \le C_p\bigl(\limsup P_k^{\rm alt}\bigr)^{1-1/p}.
\]
For $k=4$, choose $p$ large in the upper bound
$P_4^{\rm alt}\le C(r/R)^{1+\alpha}$.  For $k=6$, five of
the alternating arms and the remaining arm have disjoint site
witnesses.  BKR for the fair product law and the five-arm and
one-arm bounds give
$P_6^{\rm alt}\le C(r/R)^{2+\beta}$.
Another sufficiently large $p$ preserves an exponent strictly
larger than two.

Finally, an all-one-color $k$-arm event is a disjoint occurrence
of the event of $k-1$ such arms and the event of one such arm.
Reimer's inequality at density $1/2$ bounds its probability by
the intersection of the first event with the color-reversal of
the second \cite{reimer00}.  That intersection supplies a mixed
$k$-arm family, since opposite colors are automatically
site-disjoint.  If color words are anchored, take the finite
union over the cyclic position of the opposite-colored arm.
The mixed-color estimate therefore proves the
same bounds for monochromatic families.  The bank and shadow
argument used above is topological and also works after a
homeomorphic change of coordinates with vanishing mesh.
\end{proof}

\subsection{Boundary coordinates, trimming, and extension}

To control contacts with a boundary, we need an exponent greater
than one for three disjoint arms using only fair interior vertices.
The boundary of an unexplored domain need not be a simple curve in
the physical plane, so the estimate must be formulated in conformal
coordinates.  We first justify Cardy tests in those coordinates,
then extend separated two-arm configurations to compare two nearby
Cardy events.

Exploration produces graph disks whose physical boundaries may
have several prime ends at the same point.  We formulate precisely
the coordinate setting in which Cardy and the boundary estimates
will be used.  Let $D_n$ be graph disks, and let
$f_n:\overline{\mathbb D}\to\overline{D_n}$ be conformal in the
open disk and continuous on its closure.  Assume that $f_n$ tends
uniformly to $f$, that $f$ is conformal in $\mathbb D$, and that
its continuous boundary parametrization has no constant arc.
The limiting boundary parametrization may identify different
prime ends.  We also allow a half-plane parameterization with its
usual compactification, obtained by a M\"obius change of disk
coordinates.

\begin{lemma}[Cardy tests in a boundary chart]\label{lem:chart-cardy}
In the preceding setting the pulled-back triangulations have
vanishing mesh.  In a coordinate neighborhood of a compact open
boundary interval on which $f$ is one-to-one, Jordan tests with
fair interior paths have their ordinary Cardy limits.  Their
paths may terminate at neighbors of boundary vertices, but no
prescribed boundary colors are counted as parts of those paths.

More generally, a test whose full boundary has a nonsimple
physical projection has the same conclusion whenever each of
its two complementary crossings admits cropped target arcs
and a connecting subdisk with Jordan image, with the cropped
Cardy values approaching the two complementary values of the
original test.  The cropping is performed before target arcs
are extended for a comparison.
\end{lemma}

\begin{proof}
If a connected triangle piece had vanishing physical diameter but
coordinate diameter bounded below, a Hausdorff subsequential
limit of its preimage would be a nondegenerate continuum $K$
on which $f$ is constant.  Uniform convergence of the maps gives
this conclusion.  Such a continuum cannot enter the open disk,
where $f$ is conformal and injective.  A nondegenerate continuum
in the boundary circle contains a nonconstant arc, contradicting
the hypothesis.  This proves the mesh assertion.  Ordinary
univalent distortion transfers the bulk estimates to balls in
compact subsets of the coordinate disk.

In a patch where the indicated boundary interval projects
injectively, crop the non-target sides of a coordinate test and
trim its target arcs away from the ends.  Its connecting subdisk
has Jordan image.  Approximate its interior sides by graph
crosscuts and join them to the specified bank interval in order.
The resulting cycles have Jordan boundary convergence in the
physical plane.  Apply Lemma~\ref{lem:good-sequences}, extending
the target caps only after this cropping, and trim comparison
paths at the graph boundary.  They approach the specified bank
from inside this subdisk, even if a different bank projects to
the same physical neighborhood.  Unused independent full-plane
bits may be added for the comparison tests.  Their trimmed
witnesses use only the fair interior bits of the subdisk.

For the last assertion do this separately for the two
complementary crossings.  Their lower bounds approach the two
complementary Cardy numbers; Hex duality supplies the squeeze.
Thus the assertion invokes Cardy only for Jordan domains in the
physical plane, not for a test on an unidentified covering
surface or on the whole nonsimple projection.
\end{proof}

We require an extension construction which does not condition on
the colors exposed in finding arm endpoints.  We therefore
describe its retained event as an existence event with deterministic
allowed regions.  A \emph{trim collar} is a fixed-ratio annulus,
or half-annulus in a boundary chart, with a middle circle or
semicircle and fixed-factor room on both sides.  A proper arm
traverses the collar from one boundary component to the other;
after a trim it is stopped at its first hit from the retained
side.  Small endpoint boxes will always be selected from a fixed
finite family of box configurations on the middle curve.

\begin{lemma}[Separated trims]\label{lem:trim-arms}
For every $z>0$, one can choose a finite family of configurations
of boxes on a trim curve, with fixed positive relative sizes
and separations within each configuration, so that the following
exception has limit superior probability at most $z$.
The exception is that two distinct simple arms traversing a
wider collar approach each other near the trim curve too closely
for each retained endpoint to have its own box and a tenfold
enlargement avoided by all the other arms.  In a half-annulus
also exclude visits too close to the two endpoints of the
semicircle on the base.  The conclusion applies to two arms in
a half-annulus and to four arms, with fixed cyclic order, in a
full annulus.  It controls extra approaches of the arms to the
endpoint boxes, as well as their first trim locations.  The
exception may be chosen as a union of local one-arm and four-arm
tests in a fixed wider collar, whose width does not depend on $z$.
\end{lemma}

\begin{proof}
We first record the boundary one-arm estimate, which does not
use trimming.  In a chart of Lemma~\ref{lem:chart-cardy}, a
half-annulus crossing has Cardy's limit.  Its logarithmic image
is a rectangle of aspect ratio $\pi^{-1}\log(R/r)$.  The
rectangle cross ratio in Cardy's formula gives a long-crossing
bound $C\exp(-\pi t/3)$ at aspect ratio $t$, and hence
\begin{equation}\label{eq:boundary-one-preliminary}
 \limsup\PP(\text{one free half-plane arm from $r$ to $R$})
       \le C(r/R)^{1/3}.
\end{equation}
Margins give the same upper bound for graph cuts approximating
the half-annulus.  Only free interior vertices are used.

At either base endpoint of a trim semicircle, a visit within
relative distance $d$ by an arm continuing across the collar
requires a one-arm event from scale $d$ to a fixed scale.
Equation~\eqref{eq:boundary-one-preliminary} makes the combined
probability arbitrarily small by first choosing $d$ small.
Away from these endpoint neighborhoods, cover the trim curve
by $O(1/d')$ balls of relative radius $d'\ll d$.  If two
distinct arms approach in such a ball, their continuations in
both directions give four disjoint site arms from scale $Cd'$
to a fixed collar scale.  Lemma~\ref{lem:color-arms} bounds the
sum by $C_d(d')^{c_0}$, which tends to zero.  Distortion is
uniform on these balls because their radii are small compared
with their distance from the coordinate boundary.  In a full
annulus there are no base endpoints, so only this second
argument is needed.

Define the bad-collar event to be exactly the union of the endpoint
one-arm tests and the four-arm tests on this grid.  These tests
are all contained in a fixed enlargement of the collar; making
$d,d'$ smaller does not enlarge that region.  Choose a grid of
boxes finer than the resulting separation tolerance, and take
the finite family of all ordered selections with the required
separations.
Properly trim the simple arms between the chosen curves.  Any
approach of a different arm to an endpoint box would have
triggered the same four-arm test: that arm and the arm ending
there both continue out of the wider collar in both directions.
Thus the tenfold enlargements may be reserved for their own
arms alone.  The retained portions use sites only on the
retained side of the trim, up to a vanishing mesh error absorbed
inside these boxes.  Proper crossings preserve their linear
or cyclic order.
\end{proof}

\begin{lemma}[Local guarded extension]\label{lem:local-guarded-extension}
Fix the endpoint boxes and separations of
Lemma~\ref{lem:trim-arms}, and prescribe separated destination
intervals or patches on the other side of fixed-ratio extension
regions, in the same order.  In a full annulus a specified cyclic
offset is allowed.  Let $H$ be the event that proper retained
arms with specified colors exist between the trim curves, with
endpoints in these boxes, and with every other arm avoiding each
box's reserved enlargement.  Require that each extension region
overlap the retained region only in the corresponding reserved
enlargement, where only the arm of its prescribed color is allowed.
The arms can be extended to the destinations,
with circuits or transverse crossings at their ends available
for further connections, at a probability cost bounded below
by a positive constant:
\begin{equation}\label{eq:local-extension}
 \PP(H\text{ and the prescribed extensions})\ge c\,\PP(H).
\end{equation}
The constant depends only on the normalized box and track layouts:
it is uniform for similarity copies of a fixed finite family of
layouts, including layouts in admissible coordinate charts.  The
bound holds for all sufficiently fine meshes of a deterministic
sequence, with the mesh threshold allowed to depend on the fixed
scales and chart.  It remains valid with additional restrictions
on bits outside the extension regions.
\end{lemma}

\begin{proof}
Around each endpoint box choose two slightly larger boxes.  In
the annular gap between them require a circuit of the endpoint's
color, surrounding the endpoint box.  The retained path reaches
the small box from outside the larger one, so this circuit
intersects it.  Begin a corridor inside the circuit and require
a crossing chain from there to the prescribed destination.
This is the lasso construction: the circuit connects every
possible retained witness to the corridor, without specifying
where the witness first entered its endpoint box.

Route the corridor normally away from the trim until it is on
the extension side with clearance, and then route it along a
polygonal open track to its destination.  The preserved order
makes the tracks disjoint.  Boxes can be made small compared
with the track layout, or the tracks can taper near the boxes.
Further lassos or transverse crossings may be placed at the
destinations.  Their rectangle chains have positive lower
bounds by Lemma~\ref{lem:good-sequences}; boundary destinations
use Lemma~\ref{lem:chart-cardy}.  Cardy's values depend only on
the normalized coordinate shapes, so these lower bounds are
uniform over the stipulated layouts and do not depend on the
physical distortion of a chart.  All tracks have positive
widths.  No path is required to follow a prescribed curve
exactly.

Let $U_B,U_W$ be the sets of sites used by the black and white
extension requirements.  These sets are disjoint for fine
meshes.  Where an extension meets a possible retained arm,
it lies inside that arm's reserved enlargement; all arms of
the other color avoid it.  Consequently, after fixing the bits
outside $U_B\cup U_W$, the existence event $H$ is increasing
in black bits on $U_B$ and in white bits on $U_W$.
The extension event has exactly the same coordinate
monotonicities.  Give the white coordinates the reversed
order and apply Harris' inequality to this product space.
The extension event itself depends only on these coordinates
and has probability at least $c>0$.  Thus the conditional
probability of its intersection with $H$ is at least $c$
times the conditional probability of $H$.  Integration proves
\eqref{eq:local-extension}, also in the presence of restrictions
outside these coordinates.  A bounded finite number of choices
of endpoint boxes is handled by summing, or by retaining a
choice of maximal probability.  At no point do we condition on
an already selected arm or assert that its unexplored side has
a fresh law.
\end{proof}

\begin{lemma}[Free boundary-arm bounds]\label{lem:boundary-arms}
In the coordinate patches of Lemma~\ref{lem:chart-cardy}, let
$A_{j,n}^+(r,R)$ denote the probability of $j$ disjoint arms
across a half-annulus, of any specified color order, using
only fair interior vertices.  With fixed positive radii,
constant-factor margins, and room within the patch,
\begin{equation}\label{eq:boundary-arms}
 \begin{aligned}
 \limsup_n A_{1,n}^+(r,R)&\le C(r/R)^{1/3},&
 \limsup_n A_{2,n}^+(r,R)&\le C(r/R),\\
 \limsup_n A_{3,n}^+(r,R)&\le C(r/R)^{4/3}.&&
 \end{aligned}
\end{equation}
Base vertices of prescribed boundary color are excluded.  The
same convention applies to a crosscut at which an exploration
is stopped before querying its vertices.  In a full-plane
tiling, arms allowed to approach a straight or smooth test
boundary are covered by slightly extending the coordinate
half-plane, with the resulting errors absorbed at the inner
radius.
\end{lemma}

\begin{proof}
The one-arm assertion is
\eqref{eq:boundary-one-preliminary}.  By
Lemma~\ref{lem:rectangle-colors}, applied to a graph half-annulus
viewed as a rectangle, it suffices for two arms to consider
opposite colors.  This use is valid after deleting all base
vertices; it requires no law for them.

First suppose the two arms have separated trims at scales
comparable to $r$ and $R$, with fixed-factor extension room.
Lemma~\ref{lem:local-guarded-extension} extends their small
ends to disjoint intervals on the base within distance $Cr$
of each other, and their large ends to two prescribed separated
arcs on a boundary at scale $R$.  There are only finitely many
endpoint-box and color-order choices.  We claim that the
extended event has probability at most $Cr/R$ in the limit.

Choose a half-disk at scale $R$, with the two far arcs reserved
for the two colors.  Vary one marked point $t$ on the base
through an interval containing both near endpoints.  Define
$E_t$ to be the black crossing from its far arc to the base
target arc on the right of $t$.  The other marks stay in fixed
separated positions at scale $R$.  Up to color reversal and
reflection, the near endpoints are black then white from left
to right, whereas the far black and white arcs are on the
left and right respectively.  When $t=t_-$ lies to the left
of the short near interval, the black connection gives $E_{t_-}$.
When $t=t_+$ lies to its right, the white connection is a
complementary crossing and excludes $E_{t_+}$.  These are
nested events on the same fair bits, and $t_+-t_-=O(r)$.
Terminal-side Hex duality gives the same implication on the
graph approximation.  The ordinary Cardy limits of the two
tests follow from Lemma~\ref{lem:chart-cardy}.  Their cross
ratios stay in a compact subinterval of $(0,1)$, and Cardy's
formula is differentiable there, with derivative in the
normalized moving mark bounded by a constant.  Therefore
\[
 \limsup\PP(\text{extended two-arm event})
 \le \lim_n\bigl[\PP(E_{t_-})-\PP(E_{t_+})\bigr]
 \le C r/R.
\]
Smooth boundary pieces flat on the relevant base interval
avoid any corner issue in this comparison.  The extension
lower bound now proves the same estimate for the separated
retained event.

It remains to bound the trimming exceptions without assuming
the estimate being proved.  Let $F(k)$ be the supremum of the
limit superior of the two-arm probability over the admitted
half-annuli of ratio $2^k$, allowing graph cuts converging to
their nominal boundaries.  Fixed-factor changes of radii
alter $k$ by a bounded integer.  Choose the two end collars
with disjoint buffers.  Each bad collar has probability at
most $z$ by Lemma~\ref{lem:trim-arms}.  On that failure keep
the necessary two-arm event in the remaining half-annulus,
deleting at most $k_0$ end octaves, where $k_0$ is absolute.
The bad-collar test and the remaining arms use disjoint sites.
Thus, for sufficiently large $k$,
\begin{equation}\label{eq:boundary-bootstrap}
 F(k)\le C_z2^{-k}+2z F(k-k_0).
\end{equation}
Intermediate graph cuts converging to dyadic semicircles
justify the same recurrence when the original ends are not
exact circles.  Choose $z$ so that $2z2^{k_0}<1/2$, and then
choose $D$ large enough for the finitely many initial ratios
and for $C_z$.  Induction in \eqref{eq:boundary-bootstrap}
gives $F(k)\le D2^{-k}$.  This proves the two-arm bound for
all ratios, with the mesh limit always taken first.

Finally three disjoint paths contain disjoint witnesses for
a two-arm event and a one-arm event.  BKR and the first two
estimates give the exponent $1+1/3=4/3$.  Color order in the
rectangle makes this valid for every prescribed three-color
list.  Small changes of the base or graph boundary are
handled with margins.  For an upper bound permitting paths
right up to a discretized physical boundary, one can extend
the fair tiling slightly beyond that boundary and use the
same Jordan sandwiches; no prescribed boundary colors are
thereby counted as free arms.
\end{proof}

Two consequences complete the interface with the later sections.
First, for an oblique straight side on the triangular lattice,
Lemma~\ref{lem:boundary-arms} gives the three-arm estimate at
fixed scales for every orientation.  The estimate is uniform
over rotations and translations: otherwise choose violating
meshes, extract a convergent subsequence of angles and of the
translation modulo a fundamental cell, and sandwich its tests
between slightly enlarged tests of the limiting orientation.
The fixed-scale estimate with margins contradicts the violation.
Scale the
lattice and iterate independent buffered blocks of one fixed,
sufficiently large ratio.  Above a sufficiently large lattice
scale this yields any exponent strictly below $4/3$; omitted
bounded lattice scales are absorbed in the constant.  A bounded
shift of the line absorbs the last lattice approximation error.
Choose
\[
 5/4<g_\triangle<\gamma_\triangle<4/3.
\]
The resulting half-plane three-arm bound
$C(r/R)^{\gamma_\triangle}$ is at most $C A_e(r,R)$ by the
triangular lower bound with exponent $g_\triangle$.
For Poisson geometry rotation invariance and the exponent two
give the same comparison.  This verifies, in every orientation,
the side-probe comparison used in Lemma~\ref{lem:localization}
and later in the site expansion.  The common exponent $g<2$
in that expansion may be enlarged to cover both models; it
need not be the sharper $g_\triangle$ used in this comparison.

Second, at a fixed exceptional chart point, where the neighboring
boundary arcs are only separately injective, a free arm coming
from positive coordinate distance has probability tending to
zero as its approach neighborhood shrinks.  Indeed such a path
has physical diameter bounded below: otherwise a sequence of
its connected coordinate traces would give a nondegenerate
continuum fiber of $f$, excluded in the proof of
Lemma~\ref{lem:chart-cardy}.  Meanwhile the physical images of
the shrinking inner neighborhoods shrink by continuity of $f$.
The ordinary physical one-arm bound therefore proves the
claim.  This also treats endpoints of chart patches.  These
fixed-point estimates, the bulk bounds of
Lemma~\ref{lem:color-arms}, and the free boundary bounds of
Lemma~\ref{lem:boundary-arms} are the estimates used to rule
out degeneracies of the explorations in Section~\ref{sec:critical}.

\subsection{Pivotal normalization}

We return to the original laws of the two models.  In the Voronoi
model, the expectations in this subsection average geometry as well
as colors.  The preceding estimates also give the deterministic
normalization used when density is varied.  We first work with a rectangle of
aspect ratio two, for which the pivotal estimate is available
directly; Section~\ref{sec:identification} will transfer the resulting
normalization to the unit square in the theorem.

\begin{lemma}[Finiteness and the auxiliary pivotal scale]\label{lem:pivotal-scale}
For every mesh, the expected number of color-pivotal sites for $Q_*$
and for
\[
 Q_0=[-1,1]\times[-1/2,1/2]
\]
with left--right target sides is finite and strictly positive.  At
criticality,
\begin{equation}\label{eq:aux-pivotal-scale}
 \EE N_e(Q_0)\asymp m_e A_e(b_e,1).
\end{equation}
Consequently, with
\begin{equation}\label{eq:aux-scale-def}
 a_e=A_e(b_e,1),\qquad q_e=\frac1{\EE N_e(Q_0)},
\end{equation}
we have $q_e\asymp(m_ea_e)^{-1}$ and $q_e\to0$.
\end{lemma}

\begin{proof}
Only finitely many cells meet a given compact rectangle almost
surely.  Their number has moments bounded by a polynomial in $e^{-1}$
as $e\downarrow0$: localize in successively larger boxes using the
occupancy screens, and combine their exponential failure tails with
the ordinary Poisson count moments.  At any fixed mesh the same
argument gives finite moments.  Conditional on this finite set, the
crossing indicator is a nonconstant monotone Boolean function: it
is zero when all cells are white and one when all are black.  Some
edge of a Boolean path from the all-white state to the all-black
state is therefore pivotal.  Every coloring has positive fair-color
probability, proving strict positivity of the conditional mean and
hence of the mean.  The triangular assertion is immediate from its
finite number of relevant cells.

For Poisson--Voronoi percolation,
\cite[Lemma~6.1 and equation~(6.1)]{vanneuville-scaling} estimates
the derivative of the crossing probability of a $2:1$ rectangle by
$R^2\alpha_{4,p}(R)$ at unit intensity.  Its derivative is with
respect to the color parameter, so Russo's formula counts precisely
the color pivotals used here.  To include $p=1/2$, which is an open
endpoint in that statement, fix $R$ and let $p\downarrow1/2$.
The annealed correlation length tends to infinity for the sufficiently
small defining crossing constant used there.  Localization and the
integrable cell count justify continuity of the derivative and of
the finite-scale arm probabilities.  Rescaling now gives
\eqref{eq:aux-pivotal-scale}.  The usual rectangle event agrees with
our intersection-with-the-rectangle convention, apart from null
general-position exceptions.

On the triangular lattice, Kesten's pivotal/four-arm
method~\cite[Section~4]{kesten87}, together with arm separation and
extension as in~\cite{nolin08}, yields the same estimate.
We verify the exact boundary convention directly.  For
the lower bound, every site in a small fixed central patch can
send four separated alternating arms from distinct neighbors to
macroscopic landing intervals directed toward the four sides.
Inner arm separation, followed by a finite prescription at a
bounded lattice radius, connects to the neighbors without using
the site itself.  Outer landing rectangles and disjoint RSW tracks
continue the arms across the four side arcs away from corners.
The probability is at least $c a_e$.  A black setting makes the
black crossing; a white setting makes an open white crossing with
clearance, so this is pivotal for the stated closed-cell event.
For the upper bound, use the straight-side three-arm estimate in
every orientation established above.  Write its exponent as
$\gamma>g$, choosing the triangular $g<4/3$.  For a relevant lattice site let $u$ be
its distance to the nearest rectangle corner, capped at a small
constant and bounded below by $b_e$, and let $d\le u$ be its
distance to the nearest side, also bounded below by $b_e$.
Lemma~\ref{lem:pivotal-arms} and independent buffered probes give
\[
 \PP(\text{the site is pivotal})
 \le C A_e(b_e,d)(d/u)^\gamma
 \le C a_e d^{-g}(d/u)^\gamma.
\]
When $d\asymp u$ the half-annulus is omitted.  The number of sites
in a dyadic $(u,d)$ class is at most $C m_eud$, including the
bounded-width row just outside the rectangle whose hexagons
meet it.  Summing first over dyadic $d\le u$ gives at most
$C m_ea_eu^{2-g}$; summing over $u$ is bounded since $g<2$.
The bounded number of corner-scale sites is covered by the
same estimate and the lower bound $a_e\ge c b_e^g$.
This proves the upper estimate with the exact closed-cell convention.
Finally \eqref{eq:four-powers} gives
$m_ea_e\ge c e^{g-2}\to\infty$, proving the last assertion.
\end{proof}
